\documentclass[11pt,a4paper]{article}
\usepackage[margin=18mm,headheight=15pt]{geometry}
\usepackage[T1]{fontenc}
\usepackage{lmodern,amsmath,amssymb,booktabs,array,enumitem,fancyhdr,xcolor}
\usepackage[american]{circuitikz}
\usepackage[hidelinks]{hyperref}
\definecolor{ink}{HTML}{17324D}
\pagestyle{fancy}\fancyhf{}\lhead{\small Electronic Circuits I}\rhead{\small Practice midterm | 30 marks}\cfoot{\small\thepage}
\setlength{\parindent}{0pt}\setlength{\parskip}{6pt}
\setlist[enumerate]{leftmargin=*,itemsep=7pt,topsep=4pt}
\newcommand{\q}[2]{\par\vspace{3pt}{\large\bfseries\color{ink}#1}\hfill\textbf{[#2 marks]}\par\hrule\vspace{7pt}}
\newcommand{\m}[1]{\hfill\textbf{[#1]}}
\newcommand{\kohm}{\,\mathrm{k}\Omega}
\newcommand{\uF}{\,\mu\mathrm{F}}
\begin{document}

{\Huge\bfseries\color{ink}Electronic Circuits I}\par
{\Large Practice midterm}\hfill\textbf{90 minutes | 30 marks}

\textbf{Scope:} Lectures 1--8 and Sheets 1--7, including the additional op-amp problems, DACs and ADCs. This is an original practice paper, not an official or predicted exam.

\textbf{Name:} \rule{8cm}{.4pt}\qquad\textbf{ID:} \rule{3cm}{.4pt}

\textbf{Instructions.} Answer all five questions on separate answer sheets. Show derivations and units. A scientific calculator is allowed. Ideal op-amps are assumed unless stated otherwise; assume adequate supply rails. Treat coupling and bypass capacitors as open circuits for DC and shorts for midband AC. Use $V_T=25$ mV. State any additional approximation.

\begin{center}\begin{tabular}{lcc}\toprule
Question & Marks & Suggested time\\\midrule
1. Op-amp synthesis and limits & 5 & 15 min\\
2. DAC and ADC & 7 & 21 min\\
3. Voltage regulation & 3 & 9 min\\
4. 555 timer & 4 & 12 min\\
5. BJT bias and amplifiers & 11 & 33 min\\\bottomrule
\end{tabular}\end{center}
\q{Question 1. Op-amp synthesis and limits}{5}
\begin{enumerate}[label=(\alph*)]
\item Using \textbf{at most three ideal op-amps}, design a circuit implementing
\[
v_o(t)=(1\,\mathrm{ms})\frac{dv_1}{dt}+(100\,\mathrm{s}^{-1})\int_0^t v_2(\tau)\,d\tau-2v_3(t).
\]
Assume zero initial capacitor voltage. Draw the complete circuit, label all components, and give one valid set of resistor and capacitor values. Justify the signs of all three terms. \m{3}
\item A noninverting amplifier has $R_f=90\kohm$, $R_g=10\kohm$, and finite open-loop gain $A_0=1000$. Derive its closed-loop gain and calculate its percentage reduction relative to the ideal gain. \m{1}
\item A conventional inverting BJT log amplifier uses an input resistor and a transistor feedback element that carries positive current for $v_{\rm in}>0$. Explain why its logarithmic formula cannot be extended to a negative input voltage. \m{1}
\end{enumerate}
\vfill\textit{Use the remaining time after each question to check signs, operating regions, and loading.}
\newpage
\q{Question 2. Digital/analog conversion}{7}

\begin{center}\begin{circuitikz}[scale=.84,transform shape]
\draw (0,0) node[ground]{} to[R,l={$2R$}] (2,0) to[R,l={$R$}] (4,0) to[R,l={$R$}] (6,0) to[R,l={$R$}] (8,0);
\foreach \x/\b in {2/0,4/1,6/2,8/3}{\draw (\x,0) to[R,l={$2R$},*-o] (\x,2) node[above]{$b_{\b}V_{\rm ref}$};}
\draw (9.5,-.49) node[op amp](a){};
\draw (8,0) -- (a.-) (a.+) -- ++(0,-.7) node[ground]{};
\draw (8.6,0) to[short,*-] (8.6,2.7) to[R,l={$R_f$}] (11.5,2.7) -- (11.5,-.49) to[short,*-o] (12,-.49) node[right]{$v_o$};
\draw (a.out) -- (11.5,-.49);
\end{circuitikz}\par\small Figure 1. Four-bit R--2R DAC; $b_3$ is the MSB.\end{center}

In Figure 1, each bit $b_i$ is either 0 or 1; its input terminal is connected to $b_iV_{\rm ref}$. Assume ideal switches.
\begin{enumerate}[label=(\alph*)]
\item With $(b_3b_2b_1b_0)=0010$, derive $v_o$ using successive Th\'evenin reductions or equivalent current division. Then write the general expression for $v_o$ for any input code. \m{2}
\item Choose $R$, $2R$, $R_f$ and a \textbf{positive} $V_{\rm ref}$ so the resistance seen at the MSB terminal is $20\kohm$ and the output spans exactly $0$ to $-3.75$ V. Use $R_f=R$. State the signed one-LSB step and the output at code $1011$. \m{2}
\end{enumerate}

\begin{center}\begin{circuitikz}[scale=.70,transform shape]
\draw (0,5.4) node[above]{$4\,\mathrm V$} to[R,l={$R$}] (0,4) to[R,l={$R$}] (0,2.3) to[R,l={$R$}] (0,.6) to[R,l={$R$}] (0,-1) node[ground]{};
\draw (1.3,5.1) node[above]{$v_{\rm in}$} -- (1.3,1.1);
\foreach \y/\n in {4/3,2.3/2,.6/1}{
\draw (2,\y-.35) -- (2,\y+1) -- (3.3,\y+.325) -- cycle;
\node at (2.2,\y+.65) {$+$};\node at (2.2,\y) {$-$};
\draw (0,\y) to[short,*-] (2,\y);
\draw (1.3,\y+.65) to[short,*-] (2,\y+.65);
\draw (3.3,\y+.325) -- (5,\y+.325) node[midway,above]{$C_{\n}$};
}
\draw (5,0) rectangle (8,5.1);\node[align=center] at (6.5,2.55){Thermometer-to-\\binary encoder};
\draw[->] (8,2.55) -- (9.2,2.55) node[right]{$D_1D_0$};
\end{circuitikz}\par\small Figure 2. Two-bit flash ADC using ideal op-amps as comparators.\end{center}

The divider uses four equal resistors. The encoder outputs the binary count of comparators that are high. Define $C_k=1$ when $v_{\rm in}\geq k$ V, for $k=1,2,3$.
\begin{enumerate}[label=(\alph*),start=3]
\item Complete the table and sketch $D_1$ and $D_0$ versus sample number, holding each code until the next sample. \m{2}
\begin{center}\begin{tabular}{c*{6}{c}}\toprule
Sample $n$ &1&2&3&4&5&6\\
$v_{\rm in}$ (V)&0.4&1.2&2.6&3.4&2.2&0.8\\
$C_3C_2C_1$&&&&&&\\
$D_1D_0$&&&&&&\\\bottomrule\end{tabular}\end{center}
\item State the ADC step size and the number of comparators required for a 3-bit version. Explain why $v_+=v_-$ must \textbf{not} be imposed on these comparator op-amps. \m{1}
\end{enumerate}
\newpage
\q{Question 3. Voltage regulation}{3}
A 7812 regulator is assisted by a PNP external pass transistor. Its emitter connects to $V_{\rm in}=18$ V, its collector to the regulated output, and its base to the IC input. A resistor $R_{\rm ext}$ connects the emitter to the base. The IC ground pin is grounded. Use $V_{EB}=0.70$ V; neglect base and ground-pin currents. In the simplified current-sharing model, the IC carries 0.40 A once the pass transistor turns on.
\begin{enumerate}[label=(\alph*)]
\item State the nominal output voltage and select $R_{\rm ext}$ so the external transistor starts conducting at an IC current of 0.40 A. \m{1}
\item With $R_L=12\,\Omega$, calculate the load current and external-transistor dissipation. \m{1}
\item Explain why the IC's internal current limit does not by itself protect the external transistor. Describe a modification that limits the \emph{total} load current to 1.5 A, and calculate the sense resistor for a 0.70 V turn-on threshold. \m{1}
\end{enumerate}
\q{Question 4. 555 timer design}{4}

\begin{center}\begin{circuitikz}[scale=.78,transform shape]
\draw (0,5) node[above]{$V_{CC}=9\,\mathrm V$} to[R,l={$R_A$}] (0,3.3) to[R,l={$R_B$}] (0,1.5) to[C,l={$C=100\,\mathrm{nF}$}] (0,0) node[ground]{};
\draw (3,.5) rectangle (6,4.4);\node at (4.5,2.7){\Large 555};
\draw (0,3.3) to[short,*-] (3,3.3) node[right]{7 DISCH};
\draw (0,1.5) to[short,*-] (2,1.5) -- (2,2) -- (3,2) node[right]{6 THRESH};
\draw (2,1.5) to[short,*-] (3,1.5) node[right]{2 TRIG};
\draw (0,5) to[short,*-] (5.4,5) -- (5.4,4.4) node[below]{4};
\draw (4.2,5) to[short,*-] (4.2,4.4) node[below]{8};
\draw (4.5,.5) node[above]{1} -- (4.5,0) node[ground]{};
\draw (6,2.7) node[left]{3} to[short,-o] (7.2,2.7) node[right]{$v_o$};
\end{circuitikz}\par\small Figure 3. Standard astable connection; pin-5 decoupling is omitted.\end{center}

\begin{enumerate}[label=(\alph*)]
\item State the lower and upper capacitor switching voltages. Design $R_A$ and $R_B$ for $f=1$ kHz and a high-state duty cycle of 75\%, using $C=100$ nF. Show the timing equations. \m{2}
\item Redesign for the same frequency and capacitance but a duty cycle of 40\%. Draw the necessary modification, specify diode orientation, and calculate the new resistor values. Use an ideal diode and ideal discharge switch. Explain why the unmodified circuit cannot achieve 40\%. \m{2}
\end{enumerate}
\newpage
\q{Question 5. BJT bias and small-signal analysis}{11}
\textbf{Parts (a)--(d):} Use Figure 4. Take $\beta=100$, $V_{BE}=0.70$ V for the DC calculation, and neglect the Early effect. The input $v_i$ is supplied by an ideal AC voltage source through $C_{\rm in}$.

\begin{center}\begin{circuitikz}[scale=.88,transform shape]
\draw (5,3.3) node[npn](q) {$Q_1$};
\draw (1.8,6) -- (5,6) node[above]{$V_{CC}=8\,\mathrm V$} to[R,l={$R_C=3\kohm$}] (q.C);
\draw (1.8,6) to[R,l_={$R_1=60\kohm$}] (1.8,3.3) -- (q.B);
\draw (1.8,3.3) to[R,l_={$R_2=20\kohm$},*-] (1.8,0) node[ground]{};
\draw (-1,3.3) node[left]{$v_i$} to[C,l={$C_{\rm in}$},o-] (1.8,3.3);
\draw (q.E) -- (5,2.2) to[R,l_={$R_E=1\kohm$}] (5,0) node[ground]{};
\draw (5,2.2) to[short,*-] (6.8,2.2) to[C,l={$C_E$}] (6.8,0) node[ground]{};
\draw (q.C) to[short,*-] (7,4.07) to[C,l={$C_{\rm out}$}] (9,4.07) node[above]{$v_o$} to[R,l={$R_L=6\kohm$}] (9,0) node[ground]{};
\end{circuitikz}\par\small Figure 4. Divider-biased common-emitter amplifier.\end{center}

\begin{enumerate}[label=(\alph*)]
\item Explain why an NPN BJT with a forward-biased base--emitter junction and reverse-biased base--collector junction can carry collector current. Identify the dominant transport mechanism through the neutral base. \m{1}
\item Replace the divider by its Th\'evenin equivalent. Calculate $I_B$, $I_C$, $V_B$, $V_E$, $V_C$ and $V_{CE}$, retaining base-current loading. Verify forward-active operation using the base and collector voltages. \m{3}
\item With $C_E$ connected, draw the midband hybrid-$\pi$ equivalent circuit. Calculate $g_m$, $r_\pi$, $A_v=v_o/v_i$, $R_{\rm in}$ seen to the right of $C_{\rm in}$, and $R_{\rm out}$ with $R_L$ removed and $v_i=0$. \m{2}
\item Remove $C_E$. Derive the new gain and input resistance, retaining finite $\beta$. State whether the DC operating point changes. \m{2}
\end{enumerate}
\textbf{Parts (e)--(f) are independent of Figure 4.}
\begin{enumerate}[label=(\alph*),start=5]
\item An emitter follower has $g_m=40$ mS, $\beta=100$, $r_o=\infty$, and $R_E=1\kohm$. Its collector is at AC ground; a coupling capacitor connects a $1\kohm$ load to its emitter. An ideal bias current source feeds its base, and a signal source of resistance $R_s=3\kohm$ drives it. Find its \emph{loaded} base-to-emitter gain and input resistance, and its output resistance with the load removed and source voltage set to zero. \m{2}
\item For a common-base stage with $I_C=1$ mA, $\beta=100$, $r_o=\infty$, an ideally AC-grounded base, and effective collector load $2\kohm$, find the gain from emitter voltage to collector voltage and the resistance looking into the emitter (no external emitter resistor). \m{1}
\end{enumerate}
\vfill\begin{center}\textbf{End of paper --- Total: 30 marks}\end{center}
\end{document}
